This paper estimates how much of India's May 2026 gold import duty hike, from 6 to 15\%, passed through to domestic gold prices. Exploiting the sharp and unanticipated policy discontinuity as a quasi-experiment, this paper applies an interrupted time series (ITS) design with Newey-West HAC standard errors to a primary sample of 336 daily observations from the India Bullion and Jewellers Association benchmark price series (July 2024 -- June 2026), with 307 pre-hike and 29 post-hike trading days. The full dataset spans January 2022 through July 2026 (824 IBJA trading days across a 1,171-row weekday date spine through July 1, 2026).

The duty hike raised the domestic gold premium by Rs.\,10,124 per 10 grams, representing 84.1\% of the theoretical ceiling, with a 95\% confidence interval of 76.6 to 91.6\%. Full pass-through is formally rejected ($t = -4.17$, $p < 0.001$). Local projections confirm the adjustment was immediate on the first trading day and held on a stable plateau through approximately the fourteenth trading day, with no evidence of anticipatory pricing. An ARDL bounds test confirms cointegration between domestic and international gold prices in the pre-hike period, establishing the shift as permanent rather than transitory. Three consistency checks using counterfactual methods trained on the pre-hike window converge within 5\% of the primary estimate, and fake-date placebo tests confirm the result is specific to May 13, 2026.

This paper also documents asymmetric price transmission: the July 2024 duty cut passed through within one week, while the hike remained incomplete at 84.1\% after 29 trading days. The incomplete pass-through is consistent with supply-side dampening from front-loaded inventory imported at the prior duty rate, seasonal demand weakness, and historical smuggling elasticity.
